

\documentclass[11pt]{article}

\usepackage[spanish]{babel}
\usepackage[utf8]{inputenc}
\usepackage[T1]{fontenc}
\usepackage{amsmath, amssymb, amsthm, mathtools}
\usepackage{geometry}
\usepackage{xcolor}
\usepackage{tcolorbox}

\geometry{letterpaper, margin=2.2cm}

% ==============================================================================
% DEFINICIONES DE ENTORNOS PERSONALIZADOS PARA COMPILACIÓN
% ==============================================================================

\newenvironment{motivacion}{%
  \section*{1. Motivación}%
}{}

\newenvironment{preguntaguia}{%
  \begin{tcolorbox}[colback=cyan!5,colframe=cyan!75!black,title=Pregunta Guía]%
}{%
  \end{tcolorbox}%
}

\newenvironment{teoria}{%
  \section*{2. Definición y Teoría}%
}{}

\newenvironment{definicion}[1]{%
  \begin{tcolorbox}[colback=blue!5,colframe=blue!75!black,title=Definición: #1]%
}{%
  \end{tcolorbox}%
}

\newenvironment{teorema}[1]{%
  \begin{tcolorbox}[colback=green!5,colframe=green!75!black,title=Teorema: #1]%
}{%
  \end{tcolorbox}%
}

\newenvironment{alerta}[1]{%
  \begin{tcolorbox}[colback=yellow!10,colframe=orange!80!black,title=Alerta: #1]%
}{%
  \end{tcolorbox}%
}

\newenvironment{procesamiento}[1]{%
  \begin{tcolorbox}[colback=gray!5,colframe=gray!60!black,title=Procedimiento: #1]%
}{%
  \end{tcolorbox}%
}

\newenvironment{aplicacion}{%
  \section*{3. Aplicación y Práctica}%
}{}

\newenvironment{ejemplo}[1]{%
  \begin{tcolorbox}[colback=white,colframe=gray!40,title=Ejemplo: #1]%
}{%
  \end{tcolorbox}%
}

% Comandos para preguntas interactivas
\newenvironment{preguntaalternativas}[1]{\subsection*{Pregunta: #1}\par\vspace{0.1cm}}{\vspace{0.3cm}}
\newcommand{\opcion}[3]{\par\noindent $\bullet$ \textbf{#1} \textit{(#2)} - #3\par\vspace{0.15cm}}

\newenvironment{preguntaverdaderofalso}[1]{\subsection*{Pregunta V/F: #1}\par\vspace{0.1cm}}{\vspace{0.3cm}}
\newcommand{\verdaderofalso}[3]{\par\noindent\textbf{Respuesta correcta: #1} \\ \textit{Si marca Falso:} #2 \\ \textit{Si marca Verdadero:} #3\par\vspace{0.15cm}}

\newenvironment{preguntacasillas}[1]{\subsection*{Selección Múltiple: #1}\par\vspace{0.1cm}}{\vspace{0.3cm}}
\newcommand{\casilla}[3]{\par\noindent $\square$ \textbf{#1} \textit{(#2)} - #3\par\vspace{0.15cm}}

% Entorno Términos Pareados (3 Columnas)
\newenvironment{pareadostrescolumnas}[1]{\subsection*{Términos Pareados: #1}}{\vspace{0.3cm}}
\newcommand{\columnaI}[1]{\textbf{Columna 1:}\begin{enumerate}#1\end{enumerate}}
\newcommand{\columnaII}[1]{\textbf{Columna 2:}\begin{itemize}#1\end{itemize}}
\newcommand{\columnaIII}[1]{\textbf{Columna 3:}\begin{itemize}#1\end{itemize}}
\newcommand{\pareotres}[2]{\textbf{Cruce #1:} #2\par}

% Entornos para ejercicios
\newenvironment{ejercicios}{%
  \section*{4. Ejercicios Resueltos y Propuestos}%
}{}

\newenvironment{ejercicioresuelto}[2]{\subsection*{Ejercicio Resuelto: #1 \small{[#2]}}}{}
\newenvironment{ejerciciodemostracion}[2]{\subsection*{Demostración: #1 \small{[#2]}}}{}
\newenvironment{ejerciciopropuesto}[2]{\subsection*{Ejercicio Propuesto: #1 \small{[#2]}}}{}

\newcommand{\enunciado}[1]{\textbf{Enunciado:} #1\par\vspace{0.2cm}}
\newcommand{\solucion}[1]{\textbf{Solución:} #1\par\vspace{0.4cm}}
\newcommand{\demostracion}[1]{\textbf{Demostración:} #1\par\vspace{0.4cm}}
\newcommand{\pista}[1]{\textbf{Pista:} #1\par\vspace{0.4cm}}

% Entorno Fórmulas
\newenvironment{formulas}{\section*{5. Fórmulas Clave}\begin{description}}{\end{description}}
\newcommand{\formula}[3]{\item[#1:] $#2$ \\ \textit{#3} \vspace{0.2cm}}

% ==============================================================================
% 1. METADATOS TÉCNICOS DE DESPLIEGUE (COMPLETAR OBLIGATORIAMENTE)
% ==============================================================================
% ID_CURSO: calculo-multivariable
% NUMERO_UNIDAD: 1
% TITULO_UNIDAD: Funciones de Varias Variables
% NUMERO_CAPITULO: 1.2
% TITULO_CAPITULO: Curvas de Nivel y Grafos
% ESTADO_BLOQUEADO: false
% ESTADO_COMPLETADO: false
% ==============================================================================

\begin{document}

% ==============================================================================
% PESTAÑA 1: MOTIVACIÓN (contentMotivation)
% ==============================================================================
\begin{motivacion}
  \textbf{¿Cómo dibujamos una montaña en una hoja de papel?}
  
  Imagina que estás planeando una expedición a la cordillera. Tienes un mapa impreso completamente plano, pero necesitas saber dónde están los acantilados más empinados y dónde se encuentran las llanuras seguras para caminar. ¿Cómo puede un simple trozo de papel de dos dimensiones mostrarte el relieve tridimensional?
  
  La respuesta que encontraron los topógrafos hace siglos son las \textbf{líneas de contorno}. Cada línea en el mapa conecta todos los puntos que están exactamente a la misma altitud. Si caminaras físicamente siguiendo esa línea, no subirías ni bajarías; te mantendrías al mismo "nivel".
  
  En matemáticas, hacemos exactamente lo mismo. Tomamos un grafo tridimensional (una superficie $z = f(x,y)$) y lo "rebanamos" con planos horizontales a distintas alturas constantes ($z = k$). Al proyectar esos cortes hacia el suelo bidimensional, obtenemos un mapa de curvas de nivel que nos revela los secretos del espacio.

  \begin{preguntaguia}
    Si miras un mapa topográfico, ¿qué significa geométricamente que dos curvas de nivel estén dibujadas muy juntas en comparación a cuando están muy separadas?
  \end{preguntaguia}
\end{motivacion}


% ==============================================================================
% PESTAÑA 2: DEFINICIÓN Y TEORÍA (contentTheory)
% ==============================================================================
\begin{teoria}
  Para entender el comportamiento de un campo escalar de dos variables $f(x,y)$, necesitamos visualizarlo. Sin embargo, dibujar directamente en tres dimensiones puede ser complejo, por lo que recurrimos a dos herramientas geométricas fundamentales: el grafo y las curvas de nivel.

  % --- DEFINICIÓN 1: GRAFO DE LA FUNCIÓN ---
  \begin{definicion}{Grafo de un Campo Escalar (Superficie)}
    El \textbf{grafo} de una función de dos variables $f(x,y)$ es el conjunto de todos los puntos $(x,y,z)$ en el espacio tridimensional $\mathbb{R}^3$ para los cuales la altura $z$ es igual al valor de la función evaluada en $(x,y)$.
    
    Formalmente:
    $$ \text{Gr}(f) = \{ (x,y,z) \in \mathbb{R}^3 \colon z = f(x,y) \wedge (x,y) \in \operatorname{dom}(f) \} $$
    En gran parte de los ejemplos que veremos, este conjunto de puntos forma una \textbf{superficie} en el espacio.
  \end{definicion}

  % --- DEFINICIÓN 2: CURVAS DE NIVEL ---
  \begin{definicion}{Curvas de Nivel}
    Una \textbf{curva de nivel} de un campo escalar $f(x,y)$ es el conjunto de todos los puntos en el dominio (plano XY) que devuelven exactamente el mismo valor constante $k$.
    
    Formalmente, la curva de nivel para una constante $k$ se define como:
    $$ C_k = \{ (x,y) \in \operatorname{dom}(f) \colon f(x,y) = k \} $$
    El valor $k$ debe pertenecer obligatoriamente a la imagen $\operatorname{im}(f)$ de la función. Al dibujar varias curvas de nivel para distintos valores de $k$ en un mismo plano cartesiano, obtenemos lo que se conoce como un \textbf{mapa de contorno}.
  \end{definicion}

  % --- ALERTA CONCEPTUAL ---
  \begin{alerta}{Curva de Nivel vs. Traza}
    ¡No confundas la dimensión! Una \textbf{traza} horizontal es la intersección física de la superficie con el plano $z = k$, es decir, es un corte que vive en el espacio \textbf{3D}. 
    
    En cambio, una \textbf{curva de nivel} es la proyección de ese corte hacia abajo, aplastado contra el piso. Por lo tanto, las curvas de nivel se dibujan \textbf{siempre en el plano 2D} (el plano $XY$).
  \end{alerta}

  % --- PROCEDIMIENTO TÍPICO ---
  \begin{procesamiento}{¿Cómo construir un Mapa de Contorno?}
    Para graficar analíticamente la familia de curvas de nivel de $f(x,y)$:
    \begin{enumerate}
      \item \textbf{Verifica la imagen:} Asegúrate de elegir valores de $k \in \operatorname{im}(f)$ que la función efectivamente pueda alcanzar.
      \item \textbf{Iguala a la constante:} Plantea la ecuación algebraica $f(x,y) = k$.
      \item \textbf{Identifica la cónica:} Manipula algebraicamente la ecuación hasta reconocer una figura geométrica conocida en 2D (rectas, parábolas, circunferencias, elipses, hipérbolas).
      \item \textbf{Grafica en el plano:} Dibuja la curva en el plano $XY$ y etiquétala explícitamente con su valor de $k$ correspondiente. Repite para distintos valores de $k$ uniformemente espaciados.
    \end{enumerate}
  \end{procesamiento}

\end{teoria}


% ==============================================================================
% PESTAÑA 3: APLICACIÓN Y PRÁCTICA (contentApplication)
% ==============================================================================
\begin{aplicacion}
  A continuación, aplicaremos los conceptos de curvas de nivel y grafos mediante un ejemplo guiado, un laboratorio 3D interactivo y una serie de ejercicios interactivos diseñados para poner a prueba tu interpretación geométrica y algebraica.

  % --- LABORATORIO 3D INTERACTIVO ---
  \begin{procesamiento}{Laboratorio 3D Interactivo \colon Planos de Corte y Curvas de Nivel}
    Manipula la superficie 3D en tiempo real, desplaza el plano de corte $z=k$ y observa cómo se generan las curvas de nivel proyectadas en el plano $XY$. Puedes seleccionar diferentes superficies (dos colinas, paraboloide, silla de montar, sombrero radial):
    
    \begin{center}
      \html{<div style="position:relative; width:100%; height:580px; border-radius:10px; overflow:hidden; border:1px solid var(--border-color); box-shadow:0 4px 16px rgba(0,0,0,0.12); margin-top:12px;">
        <iframe src="animaciones/curvas_nivel/index.html" style="width:100%; height:100%; border:none;" allowfullscreen></iframe>
      </div>
      <div style="margin-top:10px; text-align:right;">
        <a href="animaciones/curvas_nivel/index.html" target="_blank" style="padding:6px 14px; background:var(--accent-color); color:white; border-radius:6px; font-weight:600; text-decoration:none; font-size:13px; display:inline-flex; align-items:center; gap:6px;">
          🔍 Abrir Laboratorio 3D en Pantalla Completa
        </a>
      </div>}
    \end{center}
  \end{procesamiento}

  % --- EJEMPLO SIMPLE ---
  \begin{ejemplo}{Construcción de Curvas de Nivel para un Paraboloide}
    Consideremos el campo escalar dado por la función:
    $$ f(x,y) = 4 - x^2 - y^2 $$
    
    \textbf{1. Determinación de Dominio e Imagen:}
    Como la expresión es polinómica, no posee restricciones algebraicas, por lo que $\operatorname{dom}(f) = \mathbb{R}^2$. Dado que $x^2 + y^2 \geq 0$, el valor máximo de la función es $4$ (alcanzado en el origen), de modo que $\operatorname{im}(f) = (-\infty, 4]$.
    
    \textbf{2. Obtención de las Curvas de Nivel:}
    Para un valor $k \in \operatorname{im}(f)$, igualamos la función a la constante $k$:
    $$ 4 - x^2 - y^2 = k \implies x^2 + y^2 = 4 - k $$
    
    \textbf{3. Análisis Geométrico según el valor de $k$:}
    \begin{itemize}
      \item \textbf{Para $k = 4$:} Obtenemos $x^2 + y^2 = 0$, lo que representa un único punto en el origen $(0,0)$.
      \item \textbf{Para $k = 3$:} Obtenemos $x^2 + y^2 = 1$, una circunferencia de radio $1$ centrada en el origen.
      \item \textbf{Para $k = 0$:} Obtenemos $x^2 + y^2 = 4$, una circunferencia de radio $2$ centrada en el origen.
      \item \textbf{Para $k > 4$:} La ecuación $x^2 + y^2 = 4 - k$ carece de solución en los números reales, por lo que $C_k = \emptyset$.
    \end{itemize}
    Geométricamente, la familia de curvas de nivel corresponde a una serie de \textbf{circunferencias concéntricas} en el plano $XY$ que se expanden a medida que el valor de altura $k$ disminuye.
  \end{ejemplo}

  % --- TIPO 1: PREGUNTA DE ALTERNATIVAS ---
  \begin{preguntaalternativas}{Identificación de Geometría de Contorno}
    ¿Qué figura geométrica representan las curvas de nivel del campo escalar $f(x,y) = y - x^2$?
    
    \opcion{Una familia de parábolas verticales desplazadas en el eje $Y$.}{Correcto}{¡Excelente! Al plantear $y - x^2 = k \implies y = x^2 + k$, vemos que cada curva de nivel es la parábola canónica $y = x^2$ trasladada verticalmente en $k$ unidades.}
    \opcion{Una familia de circunferencias concéntricas centradas en el origen.}{Incorrecto}{Recuerda que para obtener circunferencias ambas variables deben estar elevadas al cuadrado con el mismo signo ($x^2 + y^2 = r^2$).}
    \opcion{Una familia de rectas paralelas con pendiente positiva.}{Incorrecto}{La presencia del término cuadrático $x^2$ impide que la relación sea lineal.}
  \end{preguntaalternativas}

  % --- TIPO 2: PREGUNTA DE VERDADERO O FALSO ---
  \begin{preguntaverdaderofalso}{Intersección de Curvas de Nivel}
    Dos curvas de nivel $C_{k_1}$ y $C_{k_2}$ correspondientes a valores distintos $k_1 \neq k_2$ de un mismo campo escalar $f(x,y)$ nunca pueden cruzarse ni intersectarse en el plano.
    
    \verdaderofalso{V}
      {¡Correcto! Si dos curvas de nivel se cruzaran en un punto $(x_0, y_0)$, la función tendría que tomar dos valores distintos $f(x_0, y_0) = k_1 \wedge f(x_0, y_0) = k_2$ simultáneamente, lo cual contradice la definición formal de función.}
      {Incorrecto. Piensa en la definición de función: a cada punto del dominio le corresponde una y solo una salida escalar. Si se cruzaran, habría dos alturas distintas para un mismo punto $(x_0,y_0)$.}
  \end{preguntaverdaderofalso}

  % --- TIPO 3: PREGUNTA DE CASILLAS ---
  \begin{preguntacasillas}{Propiedades de Campos Racionales}
    Considera la función $f(x,y) = \frac{y}{x}$. Selecciona \textbf{todas} las afirmaciones matemáticamente correctas respecto a su dominio y curvas de nivel:
    
    \casilla{El dominio natural de la función es $\operatorname{dom}(f) = \{ (x,y) \in \mathbb{R}^2 \colon x \neq 0 \}$.}{Correcto}{Correcto. La única restricción algebraica es la división por cero, por lo que debemos excluir todo el eje $Y$.}
    \casilla{Las curvas de nivel $f(x,y) = k$ corresponden a rectas que pasan por el origen (excluyendo el punto $(0,0)$).}{Correcto}{Correcto. La ecuación $\frac{y}{x} = k \implies y = kx$ representa rectas de pendiente $k$, pero como $x \neq 0$, el origen no pertenece al dominio.}
    \casilla{La curva de nivel para $k = 0$ es todo el eje $X$.}{Incorrecto}{¡Cuidado! Aunque $y = 0 \cdot x \implies y = 0$ (el eje $X$), el punto $(0,0)$ debe ser excluido porque no pertenece al dominio.}
    \casilla{El recorrido de la función es únicamente $\operatorname{im}(f) = [0, \infty)$.}{Incorrecto}{Incorrecto. La constante $k$ puede tomar cualquier número real (positivo, negativo o cero), por lo que $\operatorname{im}(f) = \mathbb{R}$.}
  \end{preguntacasillas}

  % --- TIPO 4: TÉRMINOS PAREADOS (3 Columnas - 5 Opciones) ---
  \begin{pareadostrescolumnas}{Asociación Avanzada: Función, Geometría de $C_k$ y Restricción de $k$}
    
    % COLUMNA 1: Funciones
    \columnaI{
      \item $f(x,y) = x^2 - y^2$
      \item $f(x,y) = \sin(x - y)$
      \item $f(x,y) = \sqrt{4x^2 + y^2}$
      \item $f(x,y) = e^{y - x^2}$
      \item $f(x,y) = \frac{-1}{x^2 + y^2}$
    }
    
    % COLUMNA 2: Geometría de la Curva de Nivel - Mezclados
    \columnaII{
      \item Parábolas cóncavas hacia arriba
      \item Hipérbolas (o par de rectas secantes si $k=0$)
      \item Circunferencias concéntricas centradas en el origen
      \item Elipses concéntricas centradas en el origen
      \item Rectas paralelas de pendiente $m=1$
    }
    
    % COLUMNA 3: Rango de constantes $k \in \operatorname{im}(f)$ - Mezclados
    \columnaIII{
      \item $k \in [-1, 1]$
      \item $k > 0$
      \item $k \in \mathbb{R}$
      \item $k < 0$
      \item $k \geq 0$
    }

    % --- SOLUCIONES Y RETROALIMENTACIÓN ---
    \pareotres{1-B-III}{¡Correcto! La ecuación $x^2 - y^2 = k$ representa una familia de hipérbolas con asíntotas $y = \pm x$ (las cuales forman la curva de nivel degenerada para $k=0$). Como la diferencia de cuadrados puede tomar cualquier valor real, $k \in \mathbb{R}$.}
    \pareotres{2-E-I}{¡Excelente! Al igualar $\sin(x-y) = k$, obtenemos $x-y = \arcsin(k) \implies y = x - \arcsin(k)$, lo que define una familia de rectas paralelas con pendiente $1$. Dado el comportamiento de la función seno, esto solo es válido si $k \in [-1, 1]$.}
    \pareotres{3-D-V}{¡Muy bien! Elevando al cuadrado $\sqrt{4x^2 + y^2} = k \implies 4x^2 + y^2 = k^2$, obtenemos la ecuación canónica de elipses centradas en el origen. Debido a la raíz cuadrada de índice par, $k$ debe ser obligatoriamente no negativo ($k \geq 0$).}
    \pareotres{4-A-II}{¡Perfecto! Aplicando logaritmo natural a $e^{y - x^2} = k$, obtenemos $y - x^2 = \ln(k) \implies y = x^2 + \ln(k)$, lo que corresponde a parábolas trasladadas verticalmente. Como la función exponencial siempre es positiva, $k > 0$.}
    \pareotres{5-C-IV}{¡Magistral! Al igualar $\frac{-1}{x^2 + y^2} = k \implies x^2 + y^2 = -\frac{1}{k}$, obtenemos circunferencias. Como $x^2 + y^2$ es estrictamente positivo en su dominio (excluyendo el origen), la fracción completa es estrictamente negativa, por lo que $k < 0$.}

  \end{pareadostrescolumnas}

\end{aplicacion}


% ==============================================================================
% PESTAÑA 4: EJERCICIOS RESUELTOS Y PROPUESTOS (contentExercises)
% ==============================================================================
\begin{ejercicios}

  % --- 1. EJERCICIO RESUELTO (Analítico / Avanzado) ---
  \begin{ejercicioresuelto}{Círculos Tangentes y Completación de Cuadrados}{nivel-3}
    \enunciado{
      Determine el dominio natural, el recorrido y describa analíticamente la familia de curvas de nivel del campo escalar definido por:
      $$ f(x,y) = \frac{2y}{x^2 + y^2} $$
    }
    \solucion{
      \textbf{1. Dominio e Imagen:}
      La única restricción es que el denominador sea distinto de cero. Como la suma de cuadrados solo es cero en el origen, el dominio es todo el plano menos ese punto:
      $$ \operatorname{dom}(f) = \{ (x,y) \in \mathbb{R}^2 \colon x^2 + y^2 \neq 0 \} = \mathbb{R}^2 \setminus \{(0,0)\} $$
      Al evaluar a lo largo del eje Y (donde $x=0$), la función se reduce a $f(0,y) = \frac{2y}{y^2} = \frac{2}{y}$. Como $y$ puede ser cualquier valor real distinto de cero, la fracción $\frac{2}{y}$ puede tomar cualquier valor real, por lo tanto, $\operatorname{im}(f) = \mathbb{R}$.

      \textbf{2. Curvas de Nivel ($C_k$):}
      Planteamos la ecuación $f(x,y) = k$ para cualquier $k \in \operatorname{im}(f)$:
      $$ \frac{2y}{x^2 + y^2} = k $$
      
      Debemos separar el análisis en dos casos:
      
      \textit{Caso A ($k = 0$):}
      $$ \frac{2y}{x^2 + y^2} = 0 \implies 2y = 0 \implies y = 0 $$
      La curva de nivel $C_0$ es el eje X ($y=0$), pero recordando la restricción del dominio, debemos excluir el origen. Así, $C_0$ corresponde a todo el eje X perforado en $(0,0)$.

      \textit{Caso B ($k \neq 0$):}
      Multiplicamos por el denominador y reordenamos los términos para completar cuadrados respecto a la variable $y$:
      \[
      \begin{aligned}
        k(x^2 + y^2) &= 2y \\
        x^2 + y^2 - \frac{2}{k}y &= 0 \\
        x^2 + \left( y^2 - \frac{2}{k}y + \frac{1}{k^2} \right) &= \frac{1}{k^2} \\
        x^2 + \left( y - \frac{1}{k} \right)^2 &= \left( \frac{1}{k} \right)^2
      \end{aligned}
      \]
      
      \textbf{Conclusión:}
      Para cada $k \neq 0$, la curva de nivel $C_k$ es una \textbf{circunferencia} de radio $R = \frac{1}{|k|}$, centrada en el punto $(0, \frac{1}{k})$. Notemos que todas estas circunferencias pasan por el origen (ya que $0^2 + (0 - 1/k)^2 = 1/k^2$), pero el origen $(0,0) \notin \operatorname{dom}(f)$, por lo que es un punto excluido ("agujero") en cada una de ellas.
    }
  \end{ejercicioresuelto}

  % --- 2. EJERCICIO DE DEMOSTRACIÓN (Teórico) ---
  \begin{ejerciciodemostracion}{Curvas de Nivel de Campos Radiales}{nivel-3}
    \enunciado{
      Sea $f \colon \mathbb{R}^2 \to \mathbb{R}$ un campo escalar con simetría radial, definido como $f(x,y) = g(x^2 + y^2)$, donde $g \colon [0, \infty) \to \mathbb{R}$ es una función estrictamente creciente. Demuestre formalmente que cualquier curva de nivel no vacía $C_k$ de $f$ es una circunferencia centrada en el origen, y determine su radio.
    }
    \demostracion{
      Por definición, la curva de nivel $C_k$ es el conjunto:
      $$ C_k = \{ (x,y) \in \operatorname{dom}(f) \colon f(x,y) = k \} $$
      Sustituyendo la regla de correspondencia de $f$:
      $$ C_k = \{ (x,y) \in \mathbb{R}^2 \colon g(x^2 + y^2) = k \} $$
      
      Dado que la función $g$ es estrictamente creciente en su dominio, sabemos que $g$ es una función \textbf{inyectiva} (uno a uno). Esto garantiza que posee una función inversa bien definida $g^{-1} \colon \operatorname{im}(g) \to [0, \infty)$.
      
      Supongamos que $k \in \operatorname{im}(f)$. Entonces, podemos aplicar la inversa $g^{-1}$ a la ecuación de la curva de nivel:
      $$ g(x^2 + y^2) = k \implies x^2 + y^2 = g^{-1}(k) $$
      
      Puesto que el dominio de $g$ (y por tanto el recorrido de $g^{-1}$) es $[0, \infty)$, sabemos con certeza que $g^{-1}(k) \geq 0$.
      Definamos la constante $R^2 = g^{-1}(k)$. Entonces la ecuación se convierte en:
      $$ x^2 + y^2 = R^2 $$
      Lo cual es exactamente la ecuación analítica de una circunferencia centrada en el origen $(0,0)$ con radio $R = \sqrt{g^{-1}(k)}$. (Si $g^{-1}(k) = 0$, la curva degenera en el único punto origen). Queda así demostrado.
    }
  \end{ejerciciodemostracion}

  % --- 3. EJERCICIO PROPUESTO 1 ---
  \begin{ejerciciopropuesto}{Curvas de Nivel Exponenciales}{nivel-2}
    \enunciado{
      Determine analíticamente la familia de curvas de nivel del campo escalar $f(x,y) = e^{xy}$. Identifique la geometría de las curvas dependiendo del signo de la constante $k$.
    }
    \pista{
      Recuerda que $\operatorname{im}(f) = (0, \infty)$, por lo que debes igualar $e^{xy} = k$ asumiendo que $k > 0$. Al aplicar logaritmo natural, obtendrás $xy = \ln(k)$. 
      Analiza qué sucede si $\ln(k) > 0$ (es decir, $k > 1$), si $\ln(k) < 0$ ($0 < k < 1$) y si $\ln(k) = 0$ ($k = 1$). Deberías obtener familias de hipérbolas equiláteras y, en el caso especial, los ejes coordenados (rectas $x=0 \vee y=0$).
    }
  \end{ejerciciopropuesto}

  % --- 4. EJERCICIO PROPUESTO 2 ---
  \begin{ejerciciopropuesto}{Cónicas y Combinación Lineal}{nivel-2}
    \enunciado{
      Considere el campo escalar $h(x,y) = 4x^2 + 9y^2$. Describa detalladamente el mapa de contorno que se genera al graficar las curvas de nivel para $k \in \{0, 36, 144\}$.
    }
    \pista{
      Aquí $\operatorname{im}(h) = [0, \infty)$. Al igualar $4x^2 + 9y^2 = k$, deberás dividir toda la ecuación por $k$ para llevarla a la forma canónica $\frac{x^2}{a^2} + \frac{y^2}{b^2} = 1$. Verás que para $k>0$ el mapa de contorno corresponde a elipses centradas en el origen, donde su eje mayor descansa sobre el eje X y su eje menor sobre el eje Y. Para $k=0$, obtendrás un único punto.
    }
  \end{ejerciciopropuesto}

\end{ejercicios}


% ==============================================================================
% PESTAÑA 5: FÓRMULAS CLAVE (contentFormulas)
% ==============================================================================
\begin{formulas}

  \formula{Grafo de un Campo Escalar}
    {\operatorname{Gr}(f) = \{ (x,y,z) \in \mathbb{R}^3 \colon z = f(x,y) \wedge (x,y) \in \operatorname{dom}(f) \}}
    {Superficie tridimensional generada por la función de dos variables en el espacio.}

  \formula{Curva de Nivel}
    {C_k = \{ (x,y) \in \operatorname{dom}(f) \colon f(x,y) = k \}}
    {Conjunto de puntos en el plano $XY$ donde la función toma un valor constante $k \in \operatorname{im}(f)$.}

  \formula{Completación de Cuadrados}
    {\begin{aligned}
      x^2 + bx &= \left(x + \frac{b}{2}\right)^2 - \left(\frac{b}{2}\right)^2 \\[0.5em]
      ax^2 + bx &= a\left(x + \frac{b}{2a}\right)^2 - \frac{b^2}{4a}
    \end{aligned}}
    {Técnica algebraica indispensable para transformar expresiones cuadráticas dispersas a la forma canónica de circunferencias, elipses o parábolas desplazadas.}

  \formula{Relación Geométrica Clave}
    {\begin{aligned}
      \text{Pendiente pronunciada } &\implies \text{Curvas de nivel muy juntas} \\[0.5em]
      \text{Región plana o suave } &\implies \text{Curvas de nivel muy separadas}
    \end{aligned}}
    {Interpretación visual fundamental para conectar el mapa de contorno bidimensional con la pendiente tridimensional del grafo.}

  \formula{Invarianza por Simetría Radial}
    {f(x,y) = g(x^2 + y^2) \implies C_k \text{ son circunferencias}}
    {Cuando el campo escalar depende exclusivamente de la distancia al origen, las curvas de nivel forman circunferencias concéntricas centradas en $(0,0)$.}

\end{formulas}

\end{document}